\section{The ordinary sparse quadratic module}\label{sec:ordinary}

The ordinary box quadratic module uses $v+1$ PSD blocks for a bag of size
$v$, against up to $2^v$ for the full preordering. This section proves that
the sparse ordinary-module hierarchy still has error
$\Cf\,O_{w,d_\infty}(\log^3r/r^2)$, where $\Cf$ is the sum of the local
nonconstant Chebyshev norms (\cref{thm:mod,cor:mod-rate}). The implicit
constant and the order threshold depend only on the width and the largest
coordinate degree; they do not depend on the number of bags or the shape of
the tree. The size of the SDP grows with the number of bags, and the
absolute error through $\Cf$ can grow with it as well.

\paragraph{Why the construction of \cref{sec:kernels} does not apply.}
Positivity of a product of univariate kernels is certified in
\cref{lem:pre-density} by multiplying interval certificates, which creates
products of distinct box generators. Such products are not in the
ordinary module, and their values under an ordinary-module functional can
be negative. For example, on two variables let $L(1)=1$, let $L$ vanish on
every monomial with an odd exponent, and let
$L(x^2)=L(y^2)=L(x^4)=L(y^4)=\tfrac34$ and $L(x^2y^2)=\tfrac38$. The
moment matrix of order two is PSD: the block on $(1,x^2,y^2)$ has Schur
complement $\tfrac3{16}\bigl(\begin{smallmatrix}1&-1\\-1&1\end{smallmatrix}\bigr)$,
and the remaining diagonal entries are $\tfrac34,\tfrac34,\tfrac38$. The
localizing matrices of $1-x^2$ and $1-y^2$ on $(1,x,y)$ are
$\diag(\tfrac14,0,\tfrac38)$ and $\diag(\tfrac14,\tfrac38,0)$. So $L$ is
$\Mod_2$-positive, yet $L((1-x^2)(1-y^2))=-\tfrac18$.

Products of polynomials that are \emph{globally} sums of squares, by
contrast, are sums of squares, so a kernel that is SOS in its input variable
for every fixed output gives nonnegative densities under ordinary-module
positivity. Gribling, de Klerk and Vera~\cite{gribling2026-squared-kernels}
built such a kernel for the dense hierarchy: a squared Fejér kernel times a
geometric-series multiplier, whose mass is $1-\varrho(x)$ with a residual
$\varrho$ that is exponentially small in the length of the series. The mass
cannot be exactly one (\cref{rem:normalization-obstruction}), and in a
sparse hierarchy unequal mass defects in different bags break the exact
equality of separator marginals on which gluing relies.

\paragraph{The repair.}
We add the residual back: the kernel $\bar\Psi=\Psi+\varrho$ has mass
exactly one, so the resulting \emph{signed} densities have exactly equal
separator marginals, as in \cref{sec:kernels}. Expanding the tensor product
of $\bar\Psi$, the terms with at most one residual factor are nonnegative
under the ordinary module, since a single interval certificate times SOS
factors stays in the module. Terms with two or more residual factors may be
negative. A weighted moment Cauchy--Schwarz inequality bounds a term with
$j$ residual factors by $\delta^j\Lambda^{v-j}$, where $\delta$ is the sup
norm of the residual on the interval and $\Lambda$ bounds the Chebyshev
norm of one kernel factor. This gives a
lower bound $-\Delta_w$ on every density, with the same $\Delta_w$ in every
bag. Adding $\Delta_w$ times the reference density and renormalizing keeps
the separator marginals equal, because the reference marginals are all
equal to one. The correction costs $\Delta_w$ times the oscillation of each
local objective, which is negligible for a suitable series length.

\subsection{The squared Fejér kernel and its normalization}
\label{sec:mod-kernel}

Fix integers $s\ge2$ and $N\ge2$ with $N$ even. Let
$\chi_j=(1-j/s)_+$ for $j\ge0$ and $\bar\chi_j=1-\chi_j=\min\{j/s,1\}$.
Define
\begin{equation}\label{eq:mod-fejer}
  \Phi_s(x,u)=1+2\sum_{j=1}^{s-1}\chi_jT_j(x)T_j(u),\qquad
  M_s(x)=\int\Phi_s(x,u)^2\,d\arcs(u)=1+2\sum_{j=1}^{s-1}\chi_j^2T_j(x)^2,
\end{equation}
\begin{equation}\label{eq:mod-constants}
  C_s=1+2\sum_{j=1}^{s-1}\chi_j^2=\frac{2s^2+1}{3s},\qquad
  z_s(x)=1-\frac{M_s(x)}{C_s},\qquad
  p_{s,N}(x)=\frac1{C_s}\sum_{j=0}^Nz_s(x)^j,
\end{equation}
and
\begin{equation}\label{eq:mod-kernel}
 \begin{aligned}
  \Psi(x,u)&=p_{s,N}(x)\,\Phi_s(x,u)^2,&
  \omega(x)&=\int\Psi(x,u)\,d\arcs(u),\\
  \varrho&=1-\omega,&
  \bar\Psi(x,u)&=\Psi(x,u)+\varrho(x).
 \end{aligned}
\end{equation}
Put
\begin{equation}\label{eq:mod-D}
  D=2(s-1)(N+1),\qquad\delta=2^{-(N+1)},\qquad
  \Lambda=\frac{s^2(N+1)}{C_s}.
\end{equation}
For a univariate polynomial $q$ let $\mathcal Aq(x)=\int q(u)\Psi(x,u)\,d\arcs(u)$
and $\bar{\mathcal A}q(x)=\int q(u)\bar\Psi(x,u)\,d\arcs(u)$.

\begin{lemma}[Squared Fejér kernel]\label{lem:sos-kernel}
The following hold.
\begin{enumerate}[label=(\alph*)]
\item $\Phi_s\ge0$ on $[-1,1]^2$, and $C_s/2\le M_s\le C_s$ on $[-1,1]$;
  hence $0\le z_s\le\frac12$ on $[-1,1]$.
\item $p_{s,N}$ is a sum of squares of polynomials of degree at most
  $(s-1)N$. For each fixed $u$, $\Psi(\cdot,u)$ is a sum of squares of
  polynomials of degree at most $D/2$; its degree in $x$ is at most $D$ and
  its degree in $u$ is at most $2(s-1)$.
\item $\omega=p_{s,N}M_s=1-z_s^{N+1}$ identically and is globally SOS,
  so $\varrho=z_s^{N+1}$
  has degree at most $D$ and $0\le\varrho\le\delta$ on $[-1,1]$.
\item $\int\bar\Psi(x,u)\,d\arcs(u)=1$ identically in $x$, $\bar{\mathcal
  A}1=1$, and $\bar{\mathcal A}T_k=\mathcal AT_k$ for $k\ge1$.
\item $\chebnorm{\Psi(\cdot,u)}\le\Lambda$ for every $u\in[-1,1]$, and
  $\Lambda\ge1$.
\item For $0\le k\le s$,
  \begin{equation}\label{eq:mod-univariate}
    \chebnorm{\bar{\mathcal A}T_k-T_k}\le\eta_k:=
    \frac{N+1}{C_s}\Bigl(\frac{k^2}{s^2}+\frac{3k^2}{2s}\Bigr)
    +\sqrt{2(D+1)}\,\delta\,[k\ge1],
  \end{equation}
  where $[k\ge1]$ is one if $k\ge1$ and zero if $k=0$.
\end{enumerate}
\end{lemma}

The construction is due to Gribling, de Klerk and
Vera~\cite{gribling2026-squared-kernels}, who use it for the dense ordinary
module. We rederive its properties because the sparse argument needs them in
a precise form: $N$ must be even for the SOS identity \eqref{eq:mod-geometric}
below; SOS is claimed for $\Psi(\cdot,u)$, not for images $\mathcal AT_k$
of signed inputs, which can be negative (for $s=2$, $\mathcal AT_1(x)=
x\,p_{2,N}(x)$); and the source degree $D$ is computed directly from the
definitions.

The proof of (f) uses an exact formula for the defect of the unnormalized
squared kernel.

\begin{lemma}[Defect identity]\label{lem:mod-defect}
Let $\mathcal Kq(x)=\int q(u)\Phi_s(x,u)^2\,d\arcs(u)$. For every $k\ge0$,
\begin{equation}\label{eq:mod-defect}
  \mathcal KT_k-M_sT_k=-\bar\chi_k^2T_k
  -\sum_{l=1}^{s-1}(\bar\chi_{l+k}-\bar\chi_l)^2T_lT_{l+k}
  -\frac12\sum_{j=1}^{k-1}(\bar\chi_{k-j}-\bar\chi_j)^2T_jT_{k-j}.
\end{equation}
Consequently, for $0\le k\le s$,
$\chebnorm{\mathcal KT_k-M_sT_k}\le k^2/s^2+3k^2/(2s)$.
\end{lemma}

\begin{proof}
For $m\ge0$, orthogonality gives $\int T_m(u)\Phi_s(x,u)\,d\arcs(u)=
\chi_mT_m(x)$, with $\chi_0=1$ and $\chi_m=0$ for $m\ge s$. By
\eqref{eq:set-cheb-identities},
\[
  T_k(u)\Phi_s(x,u)=T_k(u)+\sum_{j\ge1}\chi_jT_j(x)
  \bigl(T_{j+k}(u)+T_{\abs{j-k}}(u)\bigr).
\]
Multiplying by $\Phi_s(x,u)$ and integrating in $u$ replaces each $T_m(u)$
by $\chi_mT_m(x)$:
\[
  \mathcal KT_k=\chi_kT_k+\sum_{j\ge1}\chi_jT_j
  \bigl(\chi_{j+k}T_{j+k}+\chi_{\abs{j-k}}T_{\abs{j-k}}\bigr).
\]
Also $M_sT_k=T_k+\sum_{j\ge1}\chi_j^2T_j(2T_jT_k)=T_k+\sum_{j\ge1}\chi_j^2T_j
(T_{j+k}+T_{\abs{j-k}})$. Subtracting,
\[
  \mathcal KT_k-M_sT_k=(\chi_k-1)T_k
  +\sum_{j\ge1}\chi_j(\chi_{j+k}-\chi_j)T_jT_{j+k}
  +\sum_{j\ge1}\chi_j(\chi_{\abs{j-k}}-\chi_j)T_jT_{\abs{j-k}}.
\]
For $k=0$ every term vanishes, in agreement with \eqref{eq:mod-defect}. Let
$k\ge1$ and split the last sum. Its terms with $j=l+k$, $l\ge1$, combine with
the terms $j=l$ of the middle sum to
$[\chi_l(\chi_{l+k}-\chi_l)+\chi_{l+k}(\chi_l-\chi_{l+k})]T_lT_{l+k}
=-(\chi_l-\chi_{l+k})^2T_lT_{l+k}$; these vanish for $l\ge s$. Its term
$j=k$ is $\chi_k(1-\chi_k)T_k$, which combines with $(\chi_k-1)T_k$ to
$-(1-\chi_k)^2T_k$. Its terms $1\le j\le k-1$ are
$\chi_j(\chi_{k-j}-\chi_j)T_jT_{k-j}$; symmetrizing under $j\mapsto k-j$
gives $-\frac12(\chi_j-\chi_{k-j})^2T_jT_{k-j}$. Since
$\chi_a-\chi_b=\bar\chi_b-\bar\chi_a$, this proves \eqref{eq:mod-defect}.

For the bound, every product $T_iT_j$ has Chebyshev norm one, and
$\abs{\bar\chi_a-\bar\chi_b}\le\abs{a-b}/s$. Hence the first term has norm
at most $k^2/s^2$, and the last at most $\frac12(k-1)k^2/s^2\le k^2/(2s)$
because $k\le s$. For the middle sum, $0\le\bar\chi_{l+k}-\bar\chi_l\le k/s$
and $\sum_{l\ge1}(\bar\chi_{l+k}-\bar\chi_l)=\sum_{i=1}^k(1-\bar\chi_i)\le
k$ by telescoping, so its norm is at most $(k/s)\cdot k=k^2/s$.
\end{proof}

\begin{proof}[Proof of \cref{lem:sos-kernel}]
(a) With $x=\cos\theta$ and $u=\cos\phi$,
$\Phi_s(x,u)=\frac12(F_s(\theta-\phi)+F_s(\theta+\phi))$, where
$F_s(t)=\sum_{\abs j<s}(1-\abs j/s)e^{\mathrm{i}jt}
=\frac1s\abs{\sum_{j=0}^{s-1}e^{\mathrm{i}jt}}^2\ge0$ is the Fejér kernel.
The formula for $M_s$ follows from orthogonality, and
$C_s=1+(2/s^2)\sum_{i=1}^{s-1}i^2=(2s^2+1)/(3s)$. The composed kernel
$R_s(a,c)=\int\Phi_s(a,u)\Phi_s(u,c)\,d\arcs(u)=1+2\sum_j\chi_j^2T_j(a)T_j(c)$
is nonnegative on $[-1,1]^2$. Since $2T_j(x)^2=1+T_{2j}(x)=1+T_j(T_2(x))$
and $T_j(1)=1$,
\[
  M_s(x)=\tfrac12\bigl(C_s+R_s(T_2(x),1)\bigr).
\]
For $x\in[-1,1]$ also $T_2(x)\in[-1,1]$, so $0\le R_s(T_2(x),1)\le C_s$.
This gives $C_s/2\le M_s\le C_s$ and $0\le z_s\le\frac12$.

(b) For even $N$,
\begin{equation}\label{eq:mod-geometric}
  \sum_{j=0}^Nz^j=\frac12\Bigl[1+\sum_{j=0}^{N/2-1}\bigl(z^j(1+z)\bigr)^2
  +\bigl(z^{N/2}\bigr)^2\Bigr],
\end{equation}
as one checks by collecting the even and odd powers of $z$. Substituting
$z=z_s(x)$, of degree at most $2(s-1)$, shows that $p_{s,N}$ is a sum of
squares of polynomials of degree at most $(s-1)N$. Multiplying each of these
by $\Phi_s(\cdot,u)$, of degree at most $s-1$, gives the claim for
$\Psi(\cdot,u)$.

(c) $M_s=C_s(1-z_s)$, so $p_{s,N}M_s=(1-z_s)\sum_{j=0}^Nz_s^j=1-z_s^{N+1}$.
Both $p_{s,N}$ and $M_s$ are globally SOS, hence so is their product $\omega$.
By (a), $0\le z_s^{N+1}\le2^{-(N+1)}$ on $[-1,1]$.

(d) $\int\bar\Psi(x,u)\,d\arcs(u)=\omega(x)+\varrho(x)=1$. Moreover
$\bar{\mathcal A}q=\mathcal Aq+\varrho\int q\,d\arcs$, and
$\int T_k\,d\arcs=0$ for $k\ge1$.

(e) For $\abs u\le1$, $\chebnorm{\Phi_s(\cdot,u)}\le1+2\sum_{j=1}^{s-1}
(1-j/s)=s$. Expanding $M_s$ with $2T_j^2=1+T_{2j}$ gives
$z_s=\frac{C_s-1}{2C_s}-\frac1{C_s}\sum_{j=1}^{s-1}\chi_j^2T_{2j}$, so
$\chebnorm{z_s}=(C_s-1)/C_s\le1$ and $\chebnorm{p_{s,N}}\le(N+1)/C_s$ by
\cref{lem:cheb-product}. Hence $\chebnorm{\Psi(\cdot,u)}\le(N+1)s^2/C_s=
\Lambda$. Finally $\Lambda=3s^3(N+1)/(2s^2+1)\ge1$ for $s\ge2$.

(f) For $k=0$, (d) gives $\bar{\mathcal A}T_0-T_0=0$. For $1\le k\le s$,
(d) and $\mathcal AT_k=p_{s,N}\mathcal KT_k$ give
\[
  \bar{\mathcal A}T_k-T_k=p_{s,N}\bigl(\mathcal KT_k-M_sT_k\bigr)
  +(p_{s,N}M_s-1)T_k
  =p_{s,N}\bigl(\mathcal KT_k-M_sT_k\bigr)-\varrho\,T_k.
\]
By \cref{lem:cheb-product,lem:mod-defect}, the first term has norm at most
$(N+1)C_s^{-1}(k^2/s^2+3k^2/(2s))$. For a univariate polynomial
$q=\sum_{j=0}^Dc_jT_j$, the Cauchy--Schwarz inequality and
$\int q^2\,d\arcs=c_0^2+\frac12\sum_{j\ge1}c_j^2$ give
$\chebnorm q\le\sqrt{2D+1}\,(\int q^2\,d\arcs)^{1/2}\le\sqrt{2(D+1)}\,
\norm q_\infty$. Applied to $\varrho$, with (c), the second term has norm at
most $\sqrt{2(D+1)}\,\delta$. The two terms may separately have degree above
$D$; only their sum, which is $\bar{\mathcal A}T_k-T_k$ and has degree at
most $D$ because $\bar\Psi$ has $x$-degree at most $D$ and $k\le s\le D$,
is ever evaluated by a truncated functional.
\end{proof}

\begin{remark}[Exact normalization of a globally SOS kernel forces source independence]
\label{rem:normalization-obstruction}
Let $K(x,u)$ be a polynomial such that $K(\cdot,u)$ is nonnegative on all
of $\R$ for every $u\in[-1,1]$ and $\int K(x,u)\,d\arcs(u)=1$ for all $x$.
Then $K$ does not depend on $x$. Indeed, write
$K=\sum_{l=0}^ma_l(u)x^l$ with $a_m\not\equiv0$ and suppose $m\ge1$.
Nonnegativity on $\R$ forces $m$ even and $a_m(u)\ge0$ for $u\in[-1,1]$,
while normalization forces $\int a_m\,d\arcs=0$. Since $\arcs$ has full
support and $a_m$ is continuous, $a_m=0$ on $[-1,1]$, hence identically,
a contradiction. Thus a polynomial kernel that depends nontrivially on
its input and is globally SOS there cannot have mass identically one.
The constant kernel $K=1$ is normalized and SOS, but does not approximate
nonconstant functions. This obstruction concerns globally nonnegative
source kernels; it does not rule out exactly normalized kernels certified
by an interval quadratic module. The observation is elementary and not
claimed as new; it explains why this SOS construction needs a correction.
\end{remark}

\subsection{Signed densities with exact separator consistency}
\label{sec:mod-signed}

From now on let $s\ge\max\{2,d_\infty\}$, $N\ge2$ even, and
$r\ge wD$. Then every local polynomial has total degree
$\deg f_b\le v_bd_\infty\le v_bs\le v_bD\le r$. For a bag $B$ and
$u\in[-1,1]^B$ put
\[
  \bar\Psi_B(x,u)=\prod_{i\in B}\bar\Psi(x_i,u_i),\qquad
  \bar{\mathcal A}_B=\bigotimes_{i\in B}\bar{\mathcal A},
\]
so that $\bar{\mathcal A}_Bq(x)=\int q(u)\bar\Psi_B(x,u)\,d\arcs^B(u)$. As a
polynomial in $x$, $\bar\Psi_B(\cdot,u)$ has total degree at most
$\abs BD$.

\begin{lemma}[Signed densities]\label{lem:mod-signed}
Let $(L_b)$ be a feasible ordinary-module family of order $r\ge wD$, and
define
\[
  \bar h_b(u)=L_b\bigl(\bar\Psi_{B_b}(\cdot,u)\bigr),\qquad
  \bar h^{(b)}_S(u_S)=L_b\bigl(\bar\Psi_S(\cdot,u_S)\bigr)\quad
  (S\subseteq B_b),
\]
with $\bar h^{(b)}_\varnothing=1$. Then:
\begin{enumerate}[label=(\alph*)]
\item $\int\bar h_b\,d\arcs^{B_b\setminus S}=\bar h^{(b)}_S$ for every
  $S\subseteq B_b$; in particular $\int\bar h_b\,d\arcs^{B_b}=1$;
\item $\bar h^{(b)}_{S_e}=\bar h^{(c)}_{S_e}$ for every edge $e=bc$;
\item $\bigl\lvert\int f_b\,\bar h_b\,d\arcs^{B_b}-L_b(f_b)\bigr\rvert\le
  C_b\,\Gamma_{v_b}$, where $\Gamma_v=(1+\eta)^v-1$ and
  \begin{equation}\label{eq:mod-eta}
    \eta=\frac{N+1}{C_s}\Bigl(\frac{d_\infty^2}{s^2}
    +\frac{3d_\infty^2}{2s}\Bigr)+\sqrt{2(D+1)}\,\delta .
  \end{equation}
\end{enumerate}
\end{lemma}

\begin{proof}
(a) and (b) follow exactly as in \cref{lem:pre-density}(b),(c), using
\cref{lem:sos-kernel}(d) and the fact that $\bar\Psi_{S_e}(\cdot,u)$ has
degree at most $\abs{S_e}D\le r\le2r$ in $x$.

(c) Write $f_b=c_{b,0}+g_b$ with $g_b=\sum_{\alpha\ne0}c_{b,\alpha}T_\alpha$.
By linearity, $\int f_b\bar h_b\,d\arcs^{B_b}=L_b(\bar{\mathcal
A}_{B_b}f_b)$, and $\bar{\mathcal A}_{B_b}1=1$. Since $L_b(1)=1$,
\[
  \int f_b\,\bar h_b\,d\arcs^{B_b}-L_b(f_b)
  =L_b\bigl(\bar{\mathcal A}_{B_b}g_b-g_b\bigr).
\]
Fix $\alpha$ with $\alpha_i\le d_\infty$ and put
$a_i=\bar{\mathcal A}T_{\alpha_i}(x_i)$, $t_i=T_{\alpha_i}(x_i)$. Then
$\bar{\mathcal A}_{B_b}T_\alpha-T_\alpha=\prod_ia_i-\prod_it_i=
\sum_i\bigl(\prod_{j<i}a_j\bigr)(a_i-t_i)\bigl(\prod_{j>i}t_j\bigr)$. The
Chebyshev norm is multiplicative for polynomials in distinct variables,
$\chebnorm{t_j}=1$, and $\chebnorm{a_j}\le1+\eta_{\alpha_j}$ by
\eqref{eq:mod-univariate}. Hence
$\chebnorm{\bar{\mathcal A}_{B_b}T_\alpha-T_\alpha}\le
\prod_i(1+\eta_{\alpha_i})-1\le\Gamma_{v_b}$, since
$\eta_k\le\eta$ for $k\le d_\infty$. Summing,
$\chebnorm{\bar{\mathcal A}_{B_b}g_b-g_b}\le C_b\Gamma_{v_b}$. This
polynomial has total degree at most $v_bD\le r$, so its tensor Chebyshev
expansion involves only $T_\beta$ with $\abs\beta\le r$, and
\cref{lem:cheb-moment} gives $\abs{L_b(\bar{\mathcal A}_{B_b}g_b-g_b)}\le
C_b\Gamma_{v_b}$.
\end{proof}

\subsection{Products of residuals}\label{sec:mod-residual}

The density $\bar h_b$ may be negative. Expanding the product defining
$\bar\Psi_{B_b}$ by the set $J\subseteq B_b$ of coordinates where the
residual is chosen,
\begin{equation}\label{eq:mod-expansion}
  \bar h_b(u)=\sum_{J\subseteq B_b}L_b\bigl(G_J(\cdot,u)\,H_J\bigr),\qquad
  G_J(x,u)=\prod_{i\in B_b\setminus J}\Psi(x_i,u_i),\quad
  H_J(x)=\prod_{i\in J}\varrho(x_i).
\end{equation}

\begin{lemma}[Residual products]\label{lem:residual-products}
Let $(L_b)$ be a feasible ordinary-module family of order $r\ge wD$, fix a
bag $b$ with $v=v_b$, a point $u\in[-1,1]^{B_b}$, and $J\subseteq B_b$ with
$j=\abs J$. Write $G=G_J(\cdot,u)$ and $H=H_J$.
\begin{enumerate}[label=(\alph*)]
\item If $j\le1$, then $L_b(GH)\ge0$.
\item $0\le L_b(G)\le\Lambda^{v-j}$.
\item $0\le L_b(GH^2)\le\delta^{2j}L_b(G)$.
\item $\abs{L_b(GH)}\le\delta^jL_b(G)\le\delta^j\Lambda^{v-j}$.
\end{enumerate}
\end{lemma}

\begin{proof}
By \cref{lem:sos-kernel}(b), $G=\sum_lq_l^2$ with $\deg q_l\le(v-j)D/2$.
All degrees below are collected in \cref{tab:mod-ledger}; each is at most
$r$ or $2r$ as required because $(v+j)D\le2vD\le2wD\le2r$.

(a) For $j=0$, $GH=G\in\sos_{2r}$ since $vD\le2r$. For $J=\{i\}$, the
polynomial $\varrho$ is nonnegative on $[-1,1]$ and has degree at most $D$,
which is even. \Cref{lem:interval-sos} gives $\varrho=\sigma_0+(1-x_i^2)
\sigma_1$ with $\sigma_0\in\sos_D$, $\sigma_1\in\sos_{D-2}$. Then
$GH=G\sigma_0+(1-x_i^2)G\sigma_1$ with $G\sigma_0\in\sos_{vD}$ and
$G\sigma_1\in\sos_{vD-2}$. Since $vD\le2r$, this lies in $\Mod_r(B_b)$.

(b) $L_b(G)\ge0$ by its SOS representation at the start of the proof. The polynomial $G$ has degree at most
$(v-j)D\le r$, and by \cref{lem:sos-kernel}(e) and multiplicativity of the
Chebyshev norm over distinct variables, $\chebnorm G\le\Lambda^{v-j}$.
\Cref{lem:cheb-moment} gives $L_b(G)\le\Lambda^{v-j}$.

(c) Order $J=\{i_1,\dots,i_j\}$ and write $\varrho_\ell=\varrho(x_{i_\ell})$.
Then
\[
  \delta^{2j}-H^2=\sum_{\ell=1}^j\delta^{2(j-\ell)}\,(\delta^2-\varrho_\ell^2)
  \prod_{\ell'<\ell}\varrho_{\ell'}^2 .
\]
Each $\delta^2-\varrho_\ell^2$ is nonnegative on $[-1,1]$ and has degree
at most $2D$, so \cref{lem:interval-sos} gives
$\delta^2-\varrho_\ell^2=\sigma_{0,\ell}+(1-x_{i_\ell}^2)\sigma_{1,\ell}$
with $\sigma_{0,\ell}\in\sos_{2D}$ and $\sigma_{1,\ell}\in\sos_{2D-2}$.
Multiplying the $\ell$-th term by the sum of squares
$G\prod_{\ell'<\ell}\varrho_{\ell'}^2$, of degree at most
$(v-j)D+2(\ell-1)D$, gives a sum of squares of degree at most
$(v-j+2\ell)D\le(v+j)D\le2r$ plus $(1-x_{i_\ell}^2)$ times a sum of squares
of degree at most $(v+j)D-2\le2r-2$. Only one box generator occurs in each
term. Hence $G(\delta^{2j}-H^2)\in\Mod_r(B_b)$ and
$L_b(GH^2)\le\delta^{2j}L_b(G)$. Also $GH^2=\sum_l(q_lH)^2$ with
$\deg(q_lH)\le(v-j)D/2+jD=(v+j)D/2\le r$, so $L_b(GH^2)\ge0$.

(d) For real $a,c$, $G(a+cH)^2=\sum_l(q_l(a+cH))^2$ is a sum of squares of
polynomials of degree at most $(v+j)D/2\le r$, so $L_b(G(a+cH)^2)\ge0$.
\Cref{lem:moment-cs} with weight $G$ and $V=\mathrm{span}\{1,H\}$ gives
$L_b(GH)^2\le L_b(G)L_b(GH^2)\le\delta^{2j}L_b(G)^2$ by (c). This includes
the case $L_b(G)=0$; no division is used. Part (b) completes the proof.
\end{proof}

\begin{table}[t]
\centering
\caption{Degree ledger for \cref{lem:mod-signed,lem:residual-products} in a
bag of size $v$, with $j$ residual factors. All bounds use $r\ge wD$.}
\label{tab:mod-ledger}
\begin{tabular}{@{}lll@{}}
\toprule
Polynomial in $x$ & Total degree at most & Use\\
\midrule
$\Psi(\cdot,u)$, $\bar\Psi(\cdot,u)$, $\omega$, $\varrho$ & $D$ & kernel input\\
$\bar\Psi_B(\cdot,u)$, $\bar{\mathcal A}_Bf_b$, $f_b$ & $vD\le r$ & evaluation via \cref{lem:cheb-moment}\\
$\bar\Psi_S(\cdot,u_S)$ for a separator & $\abs SD\le r$ & exact separator consistency\\
$G$ & $(v-j)D\le r$ & SOS and coefficient bound\\
$GH$ & $vD\le r$ & term of \eqref{eq:mod-expansion}\\
single-residual certificate of $GH$ & $vD\le2r$ & ordinary-module positivity\\
$GH^2$; each term of $G(\delta^{2j}-H^2)$ & $(v+j)D\le2r$ & residual bound\\
square roots $q_lH$, $q_l(a+cH)$ & $(v+j)D/2\le r$ & weighted Cauchy--Schwarz\\
\bottomrule
\end{tabular}
\end{table}

\subsection{The finite-order theorem}\label{sec:mod-theorem}

For $v\ge0$ let
\begin{equation}\label{eq:mod-Delta}
  \Delta_v=\sum_{j=2}^v\binom vj\delta^j\Lambda^{v-j},
\end{equation}
with $\Delta_0=\Delta_1=0$ (empty sums), and let
$W=\sum_b\osc(f_b)\le2\Cf$.

\begin{theorem}[Rounding for the sparse ordinary module]\label{thm:mod}
Let $s\ge\max\{2,d_\infty\}$ and $N\ge2$ be integers with $N$ even, and
define $C_s,D,\delta,\Lambda,\eta,\Gamma_v,\Delta_w$ by
\eqref{eq:mod-constants}, \eqref{eq:mod-D}, \eqref{eq:mod-eta} and
\eqref{eq:mod-Delta}. Let $r\ge wD$. For every feasible ordinary-module
family $(L_b)$ of order $r$ there is a Borel probability measure $\nu$ on
$[-1,1]^n$ whose marginal on $[-1,1]^{B_b}$ has the polynomial density
\begin{equation}\label{eq:mod-corrected}
  h_b^+=\frac{\bar h_b+\Delta_w}{1+\Delta_w}
\end{equation}
with respect to $\arcs^{B_b}$, for every bag, and
\begin{equation}\label{eq:mod-rounding}
  \Bigl\lvert\int f\,d\nu-\sum_bL_b(f_b)\Bigr\rvert
  \le\sum_bC_b\,\Gamma_{v_b}+\Delta_wW
  \le\Cf\bigl(\Gamma_w+2\Delta_w\bigr).
\end{equation}
Consequently $0\le\fstar-\rhomod_r\le\sum_bC_b\Gamma_{v_b}+\Delta_wW$.
\end{theorem}

\begin{proof}
\emph{Lower bound on the signed densities.} Fix $b$ and
$u\in[-1,1]^{B_b}$, and use \eqref{eq:mod-expansion}. By
\cref{lem:residual-products}(a) the terms with $\abs J\le1$ are
nonnegative, and by part (d) the terms with $\abs J=j\ge2$ are at least
$-\delta^j\Lambda^{v_b-j}$. Hence $\bar h_b(u)\ge-\Delta_{v_b}$. Since
$\Lambda\ge1$, the sequence $\Delta_v$ is nondecreasing in $v$; explicitly,
for $v\ge1$,
$\Delta_{v+1}=(\Lambda+\delta)\Delta_v+v\Lambda^{v-1}\delta^2\ge\Delta_v$,
and $\Delta_0=\Delta_1=0$.
So $\bar h_b\ge-\Delta_w$ on $[-1,1]^{B_b}$.

\emph{Corrected laws.} By the lower bound and \cref{lem:mod-signed}(a),
$h_b^+$ is nonnegative with integral one. Its marginal on $S\subseteq B_b$
is $(\bar h^{(b)}_S+\Delta_w)/(1+\Delta_w)$, because the constant $\Delta_w$
integrates to $\Delta_w$ over the removed coordinates. The same scalar
$\Delta_w$ is used in every bag, so \cref{lem:mod-signed}(b) gives equal
marginals on every separator, including empty ones and bags of different
sizes. For an empty bag, $\bar h_b=1$ and $h_b^+=1$, so its constant
objective is preserved exactly. \Cref{lem:gluing} (or \cref{rem:ker-gluing}) gives a probability
measure $\nu$ on $[-1,1]^n$ with marginals $h_b^+\,d\arcs^{B_b}$. A
bagwise choice of different shifts would in general destroy the equality of
separator marginals.

\emph{Objective.} Write $I_b=\int f_b\,d\arcs^{B_b}$,
$\bar I_b=\int f_b\bar h_b\,d\arcs^{B_b}$ and
$I^+_b=\int f_bh_b^+\,d\arcs^{B_b}=(\bar I_b+\Delta_wI_b)/(1+\Delta_w)$.
Then
\[
  I^+_b-\bar I_b=\Delta_w\bigl(I_b-I^+_b\bigr).
\]
Both $I_b$ and $I^+_b$ are expectations of $f_b$ under probability measures
on the bag box, so $\abs{I_b-I^+_b}\le\osc(f_b)$. No probability
inequality is applied to the signed density $\bar h_b$. With
\cref{lem:mod-signed}(c),
\[
  \Bigl\lvert\int f\,d\nu-\sum_bL_b(f_b)\Bigr\rvert
  \le\sum_b\bigl(\abs{I^+_b-\bar I_b}+\abs{\bar I_b-L_b(f_b)}\bigr)
  \le\sum_b\bigl(\Delta_w\osc(f_b)+C_b\Gamma_{v_b}\bigr).
\]
The second inequality in \eqref{eq:mod-rounding} uses $\Gamma_{v_b}\le
\Gamma_w$ and $W\le2\Cf$. The gap bound follows as in \cref{thm:pre}:
$\fstar\le\int f\,d\nu$ for every feasible family, and
\cref{lem:weak-duality} gives $\rhomod_r\le\fstar$.
\end{proof}

Additive constants in the $f_b$ cancel in every step, and no reciprocal of
a possibly small mass appears. The bound has no term that grows with the
number of edges or with the diameter of the tree: the positivity correction
couples bags only through the common scalar $\Delta_w$.

\subsection{The rate and the certificate}\label{sec:mod-rate}

The two error terms in \cref{thm:mod} pull the parameters in different
directions. The smoothing term $\Gamma_w\approx w\eta$ is of order
$N/s^2$, while the correction $\Delta_w$ is of order
$\delta^2\Lambda^{w-2}$ with $\Lambda$ of order $sN$. A series length
$N$ of order $\log s$, with a width-dependent constant, makes the
correction negligible; the degree constraint $wD\le r$ then allows $s$ of
order $r/(w\log r)$.

\begin{corollary}[Rate for the sparse ordinary module]\label{cor:mod-rate}
Let $c_w=\max\{3,(w+1)/2\}$ and $\ell_r=\log_2(r+2)$. Suppose
\begin{equation}\label{eq:mod-threshold}
  r\ge2w(c_w+3)\,\ell_r\max\{2,d_\infty\}
  \quad\text{and}\quad
  16\,w^3(3d_\infty^2+1)(c_w+3)^3\ell_r^3\le r^2 .
\end{equation}
Then
\begin{equation}\label{eq:mod-rate}
  0\le\fstar-\rhomod_r\le\Cf\Bigl[
  \frac{16e\,w^3(3d_\infty^2+1)(c_w+3)^3\ell_r^3}{r^2}
  +\binom w2\frac{2^{w+3}w^3(c_w+3)^{w+1}\ell_r^{w+1}}{r^3}\Bigr].
\end{equation}
In particular $\fstar-\rhomod_r\le\Cf\,O_{w,d_\infty}(\log^3r/r^2)$ for all
sufficiently large $r$, where the implicit constant and the threshold
depend only on $w$ and $d_\infty$. If $\Cf=0$, every $f_b$ is constant and
the gap is zero.
\end{corollary}

\begin{proof}
Let $s=\lfloor r/(2w(c_w+3)\ell_r)\rfloor$ and
$N=2\lceil(c_w/2)\log_2s\rceil$. The first condition in
\eqref{eq:mod-threshold} gives $s\ge\max\{2,d_\infty\}$. Then $N$ is even,
$N\ge2$, and $c_w\log_2s\le N<c_w\log_2s+2$. We check the hypotheses of
\cref{thm:mod} and bound its right-hand side.

\emph{Degree.} $N+1\le c_w\log_2s+3\le(c_w+3)\ell_r$, because
$\log_2s\le\ell_r$ and $\ell_r\ge1$. Hence
$wD\le2ws(N+1)\le2w(c_w+3)\ell_rs\le r$.

\emph{Smoothing term.} Since $N+1\ge c_w\log_2s+1$, we have
$\delta\le\frac12s^{-c_w}\le\frac12s^{-3}$. Using $C_s\ge2s/3$ and $s\ge2$,
\[
  \frac{N+1}{C_s}\Bigl(\frac{d_\infty^2}{s^2}+\frac{3d_\infty^2}{2s}\Bigr)
  \le(N+1)d_\infty^2\Bigl(\frac3{2s^3}+\frac9{4s^2}\Bigr)
  \le\frac{3d_\infty^2(N+1)}{s^2}.
\]
Also $D+1\le2s(N+1)$, so
$\sqrt{2(D+1)}\,\delta\le2\sqrt{s(N+1)}\cdot\frac1{2s^3}=
\sqrt{N+1}\,s^{-5/2}\le(N+1)/s^2$. Thus
$\eta\le(3d_\infty^2+1)(N+1)/s^2$. Since $r/(2w(c_w+3)\ell_r)\ge1$, the
floor satisfies $s\ge r/(4w(c_w+3)\ell_r)$, so
\[
  \frac{N+1}{s^2}\le(c_w+3)\ell_r\cdot\frac{16w^2(c_w+3)^2\ell_r^2}{r^2}
  =\frac{16w^2(c_w+3)^3\ell_r^3}{r^2}.
\]
By the second condition in \eqref{eq:mod-threshold}, $w\eta\le1$, and then
$\Gamma_w=(1+\eta)^w-1\le e^{w\eta}-1\le w\eta e^{w\eta}\le e\,w\eta$. This
gives the first term of \eqref{eq:mod-rate}.

\emph{Correction term.} For $w=1$, $\Delta_w=0$. For $w\ge2$, Taylor's
formula for $F(t)=(\Lambda+t)^w$ gives
$\Delta_w=F(\delta)-F(0)-\delta F'(0)\le\binom w2\delta^2(\Lambda+\delta)^{w-2}$.
Moreover $\Lambda\le3s(N+1)/2$, so $\Lambda+\delta\le2s(N+1)$, and
$\delta^2\le\frac14s^{-2c_w}$ with $2c_w\ge\max\{6,w+1\}$. Hence
\[
  \Delta_w\le\binom w2\frac{2^{w-2}(N+1)^{w-2}s^{w-2}}{4s^{2c_w}}
  \le\binom w2\,2^{w-4}\,\frac{(N+1)^{w-2}}{s^3},
\]
because $w-2-2c_w\le-3$. With $N+1\le(c_w+3)\ell_r$ and
$s^{-3}\le64w^3(c_w+3)^3\ell_r^3/r^3$, the term $2\Delta_w\Cf$ is at most
the second term of \eqref{eq:mod-rate}. Inserting both bounds into
$\Cf(\Gamma_w+2\Delta_w)$ proves \eqref{eq:mod-rate}. The second term is
$O_w(\ell_r^{w+1}/r^3)$, which is $o(\ell_r^3/r^2)$ for fixed $w$.
\end{proof}

The corollary states the dependence on the problem only through $\Cf$, $w$
and $d_\infty$. For bounded width and bounded coordinate degree, the order
needed to reach a prescribed error per unit of $\Cf$ is independent of the
number of bags and of the tree. For $w\le5$ the series length is
$N=2\lceil\frac32\log_2s\rceil$; larger widths need a longer series so that
$\Delta_w$ stays negligible. The constants and the threshold depend
strongly on $w$, and the corollary makes no statement for growing width.

\begin{remark}[Size of the constants]\label{rem:calibration}
The finite bound $\Cf(\Gamma_w+2\Delta_w)$ in \eqref{eq:mod-rounding} is
conservative. For $d_\infty=2$, the first three rows below use the
prescription for $N$ in the proof of \cref{cor:mod-rate} and the largest
admissible $s$ with $wD\le r$ under that prescription. The final row uses
a different admissible pair $(s,N)$. The displayed values of
$\Gamma_w+2\Delta_w$ are rounded evaluations of the worst-case formula,
not observed relaxation gaps.
\begin{center}
\begin{tabular}{@{}rrrrr@{}}
\toprule
$w$ & $r$ & $s$ & $N$ & $\Gamma_w+2\Delta_w$\\
\midrule
2 & 10{,}000 & 109 & 22 & $0.0354$\\
5 & 10{,}000 & 53 & 18 & $0.600$\\
10 & 100{,}000 & 122 & 40 & $1.87\times10^{8}$\\
10 & 100{,}000 & 82 & 60 & $1.21$\\
\bottomrule
\end{tabular}
\end{center}
The last two rows show that the prescription in \cref{cor:mod-rate} is
chosen for the asymptotic statement and is far from optimal at finite
order: for $w=10$ the admissible choice $s=82$, $N=60$ reduces the bound by
eight orders of magnitude. The orders at which the bound becomes small are
far beyond the orders used in practice; the theorem is a statement about
the rate, not a practical accuracy guarantee.
\end{remark}

\begin{corollary}[Sparse ordinary-module certificate]\label{cor:mod-certificate}
Under the hypotheses of \cref{thm:mod}, with
$E=\sum_bC_b\Gamma_{v_b}+\Delta_wW$,
\[
  f-\fstar+E\in\sum_{b=1}^t\Mod_r(B_b),
\]
with real coefficients and every summand of total degree at most $2r$.
Under the hypotheses of \cref{cor:mod-rate}, the same holds with $E$ the
right-hand side of \eqref{eq:mod-rate}.
\end{corollary}

\begin{proof}
By \cref{lem:slater}, $f-\rhomod_r\in\sum_b\Mod_r(B_b)$, and
$\rhomod_r-\fstar+E\ge0$ by \cref{thm:mod}; add this constant to the
square part of one bag. For the second statement, the proof of
\cref{cor:mod-rate} shows that $E$ is at most the right-hand side of
\eqref{eq:mod-rate} for the parameters chosen there.
\end{proof}

\subsection{Relation to prior work and scope}\label{sec:mod-prior}

\begin{remark}[What is prior in \cref{thm:mod}]\label{rem:mod-prior}
The squared Fejér kernel, its geometric-series SOS normalization, the
coefficient-norm method, and the dense rate $O(\log^3r/r^2)$ for the
ordinary box module are due to Gribling, de Klerk and
Vera~\cite{gribling2026-squared-kernels}; their theorem concerns the dense
hierarchy. The earlier dense ordinary-module rate was
$O(1/r)$~\cite{baldi2024-putinar-hypercube}. Korda, Magron and
Ríos-Zertuche~\cite{korda2025-convergence-rates-sparsity}, Theorem~8, treat
sparse ordinary modules on general semialgebraic domains under running
intersection, normalized local Archimedean certificates and local
Łojasiewicz inequalities; for the box their rate exponent is a small
negative power of $r$ that decreases with $w$. Their theorem covers
broader domains than ours. Gamertsfelder and
Mourrain~\cite{gamertsfelder2025-countable-gmp} give a general transfer of
positivity rates to generalized moment problems, which include separator
equalities, under the assumption that a dual with finitely supported
multipliers is attained; \cref{thm:mod} needs no such assumption, and its
finite SDP duality (\cref{lem:slater}) is a different statement. Magron's
lecture slides~\cite{magron2026slides-tenors} mention a sparse Putinar
rate without stating its exponent.

To our knowledge the following are new: the exactly normalized signed
kernel together with the residual-product bound of
\cref{lem:residual-products}, which uses only ordinary-module positivity,
and the resulting transfer of the dense rate to the sparse ordinary module
for arbitrary feasible truncated functionals, with an error normalized by
$\Cf$ and no dependence on the number of bags or the tree. The common
affine correction in the last step is elementary.
\end{remark}

The theorem concerns minimization over the box. The corrected laws are
supported on the box but need not satisfy additional constraints even if
the relaxation contains their localizing matrices; constrained problems are
treated in \cref{sec:extensions} under additional hypotheses. We do not
know whether the logarithmic factor in \cref{cor:mod-rate} is necessary:
the example of \cref{sec:sharpness} gives a lower bound of order $r^{-2}$
for the ordinary module, which leaves a gap of a factor $\log^3r$.
